% !TEX program = lualatex
% ECON 201, Writing Assignment Part 1: the proposal for an economic naturalist
% essay (Robert Frank, Journal of Economic Education, 2006). Compile with
% LuaLaTeX so the PDF is tagged. Build from the course root:
%   cd writing && lualatex writing1.tex && lualatex writing1.tex
% then delete the aux files and the -luamml-mathml.html side file.
% Draft: not linked from the site or the schedule yet.
\DocumentMetadata{
  pdfstandard=UA-2,
  pdfversion=2.0,
  lang=en-US,
  tagging=on
}
\documentclass[11pt]{article}

\usepackage[letterpaper, margin=0.9in, top=0.75in, bottom=0.9in]{geometry}
\usepackage{fontspec}
\setmainfont{Fira Sans}
\usepackage{microtype}
\usepackage{booktabs}
\usepackage{array}
\usepackage{enumitem}
\usepackage{xcolor}
\definecolor{accent}{HTML}{BF5700}
\definecolor{accentdark}{HTML}{8F4100}
\definecolor{ink}{HTML}{1A1A1A}
\definecolor{muted}{HTML}{595A5B}
\definecolor{cream}{HTML}{FDF3EA}
\definecolor{rulegrey}{HTML}{D9D9D9}
\usepackage[most]{tcolorbox}
\usepackage{titlesec}
\usepackage{fancyhdr}
\usepackage{lastpage}
\usepackage{hyperref}
\hypersetup{colorlinks=true, urlcolor=accentdark, linkcolor=ink,
  pdftitle={ECON 201 Writing Assignment, Part 1: Proposal},
  pdfauthor={Div Bhagia}}

\linespread{1.05}
\setlength{\parindent}{0pt}
\setlength{\parskip}{5pt}
\AtBeginDocument{\color{ink}}
\setlist{itemsep=2pt, topsep=3pt, parsep=0pt}

\titleformat{\section}{\large\bfseries\color{accent}}{}{0em}{}%
  [{\color{rulegrey}\titlerule[0.6pt]}]
\titlespacing*{\section}{0pt}{12pt plus 2pt}{5pt}

\pagestyle{fancy}
\fancyhf{}
\renewcommand{\headrulewidth}{0pt}
\fancyfoot[L]{\small\color{muted}ECON 201 \textperiodcentered\ Writing Assignment, Part 1}
\fancyfoot[R]{\small\color{muted}Page \thepage\ of \pageref{LastPage}}

% Cream panel; text inside stays ink, the title sits on the brand orange in white.
\newtcolorbox{panel}[1][]{
  enhanced, colback=cream, colframe=accent, boxrule=0.6pt, arc=2pt,
  left=8pt, right=8pt, top=6pt, bottom=6pt,
  before skip=8pt, after skip=10pt,
  title={#1}, fonttitle=\bfseries, coltitle=white,
  attach boxed title to top left={yshift=-2mm, xshift=4mm},
  boxed title style={colback=accentdark, colframe=accentdark, arc=2pt, boxrule=0pt,
    left=5pt, right=5pt, top=2pt, bottom=2pt},
  colupper=ink}

\newcolumntype{L}[1]{>{\raggedright\arraybackslash}p{#1}}
\renewcommand{\arraystretch}{1.3}

\begin{document}

{\LARGE\bfseries\color{accent}Writing Assignment, Part 1: Proposal}\hfill{\large\color{muted}ECON 201}\par
\vspace{4pt}{\color{accent}\hrule height 1.2pt}\par
\vspace{6pt}{\color{muted}Due Monday, November 16, before class, on Canvas. Worth 5\% of your course grade.}\par
\vspace{4pt}

\section{Overview}

As a part of this class, you will write a short essay as an \textbf{economic naturalist}. You will notice something puzzling in your everyday life, pose a question about it, and answer it using an idea from this course. The assignment has two parts:

\begin{itemize}
  \item Monday, November 16: Part 1, Proposal (5\%)
  \item Monday, December 7: Part 2, Essay (10\%)
\end{itemize}

The goal of the first part is to make you pick your question. I will give you feedback on your proposal, and you will use it to write your essay. The essay will be 300 to 400 words. Instructions for Part 2 will be posted on Canvas when you receive your feedback.

The idea comes from the economist Robert Frank. Someone who has taken a biology course starts noticing things in nature they never saw before. Frank wants economics to do the same for you. Here are two questions his students asked:

\begin{itemize}
  \item \textit{Why do the keypads on drive-up ATMs have Braille dots?} Banks need Braille keypads for walk-up machines anyway. Making one kind of keypad for every machine is cheaper than making two, and the dots do not bother sighted drivers.
  \item \textit{Why do brides buy expensive wedding dresses while grooms rent tuxedos?} Most brides want a distinctive dress, so a rental shop would need dozens of styles in every size, each rented rarely. It would have to charge more than the price of the dress. Grooms are happy with a standard tuxedo, so renting one is cheap.
\end{itemize}

Neither answer needs a graph or an equation. Each uses one simple idea, applied carefully to something the student actually saw.

\section{Choosing Your Question}

\textbf{Your question should be about economics.} Economics explains the choices people and businesses make about prices, quantities, and who gets what. A good question asks why someone makes a choice that looks odd once you think about costs and benefits. For example, why does a business charge this price or offer this deal? Why do people buy, sell, wait in line, or switch? Why is something shared, rationed, given away, or not sold at all?

\textbf{Your question should pass three tests.}

\begin{enumerate}
  \item \textit{It is about a pattern, not a one-time event.} You have seen it more than once, or it happens in many places.
  \item \textit{An idea from this course can answer it.} You should not need inside information about a company, such as its costs or what its managers were thinking.
  \item \textit{The answer is not obvious.} A friend who has never taken economics would be stumped or would guess wrong.
\end{enumerate}

Questions about the whole economy, inflation, interest rates, or the stock market will not work. This course is about individual markets. The table below shows some questions that do not work and how to turn them into ones that do.

\begin{center}
\small
\begin{tabular}{@{}L{4.6cm}L{4.2cm}L{6.6cm}@{}}
\toprule
\textbf{Does not work} & \textbf{Why} & \textbf{Works} \\
\midrule
Why is everything so expensive? & It is about the whole economy. & Why did egg prices double while milk prices barely moved? \\
Why do people like Labubu? & It is about taste. & Why do Labubu sell out in minutes and then resell for triple the price? \\
Should college be free? & It asks for an opinion. & Why does the campus health center give flu shots for free? \\
Why did my manager cut my hours? & Only your manager knows. & Why do stores hire extra workers in December and let them go in January? \\
\bottomrule
\end{tabular}
\end{center}

\textbf{Your question should be uncommon.} The best questions come from a place you know well and others may not, such as your job, your family's business, your neighborhood, or a hobby. Some topics come up every semester and will score lower on originality: parking, textbook prices, movie theater popcorn, printer ink, airport food, Uber surge pricing, and boba cup sizes. You may not use any question in this document.

\section{Ideas You Can Use}

Your proposal should name the course idea you think answers your question. It can come from a lecture we have not reached yet, so the later lectures are listed too. If you are not sure which idea fits, make your best guess. My feedback will point you to the right one.

\begin{center}
\small
\begin{tabular}{@{}L{3.6cm}L{12.0cm}@{}}
\toprule
\textbf{Lectures} & \textbf{Ideas} \\
\midrule
3 and 4 & Opportunity cost, economic rent, incentives, comparative advantage, gains from trade \\
5 to 9 & Product differentiation, fixed and marginal cost, economies of scale, price elasticity of demand, substitutes and complements, choosing a price to maximize profit \\
10 to 13 & Consumer and producer surplus, deadweight loss, market power, price discrimination, natural monopoly, price wars and competition \\
15 to 20 & Supply and demand, price-taking firms, who gets what, shifts in supply and demand, entry and exit, taxes and price controls \\
21 and 22 & External effects, property rights and bargaining, regulation and taxes that fix external effects \\
23 and 24 & Public goods, free riding, shared resources, hidden actions, hidden attributes, the limits of markets \\
\bottomrule
\end{tabular}
\end{center}

\section{Instructions}

Write \textbf{one paragraph of about 150 words} that does the following, in this order:

\begin{enumerate}
  \item \textbf{Describe the puzzle.} Say what you saw, where, and when. Then explain why it is puzzling: what would you expect to happen, and why does that not happen?
  \item \textbf{State your question} in one sentence.
  \item \textbf{Give your first guess at the answer} and name the course idea behind it.
\end{enumerate}

You do not need the right answer yet. The point of the proposal is to find a good question. Write in full sentences, not bullet points.

\begin{panel}[Example proposal]
On my drive to campus, I pass four car washes, and every one advertises an unlimited monthly plan for about \$20, the price of one or two washes. That seems odd. A gym can offer unlimited visits because one more visit costs it almost nothing, but every car wash uses water, soap, and electricity. If members washed their cars every few days, the plan would lose money. Yet these plans keep spreading. My question is: why would a car wash sell unlimited washes for the price of two? My first guess is that the cost of one more wash is small compared with the cost of the machines and the lot, which the owner pays either way, and that most members wash far less often than they expected when they signed up. The ideas I would use are marginal cost and fixed cost.
\end{panel}

\section{What Needs to Be Submitted}

\begin{itemize}
  \item Your paragraph, in Word or PDF format, with your name at the top.
  \item Below the paragraph, one or two sentences describing how you used AI tools, if at all. An honest description of substantial use is not penalized. An undisclosed one is a violation of academic integrity.
\end{itemize}

Submissions need to be made on Canvas. Work submitted after the due date loses 10\% per day, up to a maximum of 50\%.

\section{Grading}

\begin{center}
\small
\begin{tabular}{@{}L{12.4cm}r@{}}
\toprule
\textbf{Criterion} & \textbf{Points} \\
\midrule
\textbf{Original.} An uncommon question drawn from your own life. & 2 \\
\textbf{Puzzle explained.} It says clearly what you would expect and why that does not happen. & 1 \\
\textbf{Economic and answerable.} The question passes the three tests. & 1 \\
\textbf{Well structured.} One paragraph, in the order above, about 150 words. & 1 \\
\midrule
\textbf{Total} & 5 \\
\bottomrule
\end{tabular}
\end{center}

\end{document}
